\section{Cook--Levin: the argument in the formalization}
\label{sec:cook-levin}

This sample translates the route through the tableau, correctness, reduction and
polynomial-time files to the final completeness theorem at revision
\texttt{263e7321}. Its mathematical sources are the tableau presentations of
Arora and Barak~\cite[Section 2.3]{AroraBarak09} and
Papadimitriou~\cite[Theorem 8.2]{Papadimitriou94}.
The underlying local transition lemma and compositional efficiency rules are used as
named prerequisites below; their internal proofs are not translated here.

\begin{theorem}[Formalized Cook--Levin theorem]
SAT, with CNF encoding $\mathrm{encCNF}$, is NP-complete under the library's tidy
polynomial-time many-one reductions. UNSAT is coNP-complete under those reductions.
\end{theorem}

\subsection{The tableau and the statement it must satisfy}

Fix a deterministic binary multi-tape machine $M$ with finite control, a word $u$, a certificate length
$L$, and a time bound $T$. The machine has a read-only input tape, finitely many
two-way work tapes and a write-only output tape. Write $C_t(w)$ for its configuration
after $t$ steps on input $u\mathbin{+\!+}w$.
The tableau formula $F(M,u,L,T)$ describes times $0,\ldots,T$.
Its variables are $X_{t,\phi}$ for machine features $\phi$, and $W_j$ for certificate bits.
The basic features record the control state, input-head position, work-head positions,
work-cell symbols and a view of the output. Work positions are restricted to
$[-T,T]$: a head starting at zero moves at most one cell per step.
Auxiliary features record the symbols read and the current \emph{situation}, meaning
the active control state together with the symbols under the heads.

The formula is the conjunction of the following clause families:
\begin{enumerate}
\item Initial clauses set the start state, input position $1$, work positions $0$,
      blank work cells, and empty output.
\item At every time, one-hot clauses require exactly one feature in each basic group:
      one state, one head position, one symbol per represented cell, and one output view.
\item Before time $T$, read clauses force the actual symbols under the heads to be
      recorded. For the input tape they consult $u$, the variables $W_j$, or an end marker.
      Situation clauses force the active state and read symbols to give their situation.
\item Step clauses encode implications from terms describing the current configuration
      to basic features of the next one.
\item Final clauses require halted state and the output view corresponding to $[\mathrm{true}]$.
\end{enumerate}
These are the constructors in
\source{Complexity/Computability/Complexity/CookLevin/Tableau.lean}{\lean{tableau}}.
The intended conclusion is
\begin{equation}\label{eq:tableau}
F(M,u,L,T)\text{ is satisfiable}
\quad\Longleftrightarrow\quad
\exists w\in\{0,1\}^{L},\ C_T(w)\text{ is halted with output }[\mathrm{true}].
\end{equation}

\subsection{Why the one-hot constraints suffice}

For an assignment $\sigma$, recover $w$ by taking $\sigma(W_j)$ for $j<L$.
Define $\mathrm{Agrees}(\sigma,t,C)$ to mean that for every basic feature $\phi$,
\[
 \sigma(X_{t,\phi})=\mathrm{true}\quad\Longleftrightarrow\quad \phi\text{ holds of }C.
\]
Each basic group has exactly one feature that holds of an actual configuration,
provided its heads lie in the represented window. Suppose the assignment satisfies
the one-hot clauses and assigns true to every basic feature that holds of $C$.
In a group, let $\psi$ be the unique feature holding of $C$. Its variable is true;
exactly-one forces every other variable in the group false. Hence the assignment
agrees with $C$ on the entire group. Applying this to every group proves
\source{Complexity/Computability/Complexity/CookLevin/Correctness.lean}{\lean{Agrees.of_oneHot}}.
This is why the later induction only needs to force the features that really occur.

\subsection{From a satisfying assignment to an accepting run}

Assume $\sigma$ satisfies the tableau. We prove by induction that it agrees with
$C_t(w)$ for every $t\leq T$. The work-head bound holds at all these times because
$|\mathrm{head}_i(C_t(w))|\leq t\leq T$.
At time zero, the initial clauses set every basic feature that holds of the initial
configuration. The one-hot argument therefore gives agreement at time zero.

For the induction step, suppose there is agreement at time $t<T$.
The read clauses force the symbols actually read in $C_t(w)$: on a work tape,
use its true head-position and cell-symbol features; on the input tape, use its
true input position and the reconstructed certificate. The situation clauses then
force the actual active situation when one holds.

Now let $\phi$ be a basic feature holding of $C_{t+1}(w)$.
The local machine lemma
\source{Complexity/Computability/Machines/Turing/MultiTape/Feature.lean}{
\lean{Feature.holds_step_iff_exists_stepTerms}}
says that some conjunction in $\mathrm{stepTerms}(M,\phi)$ holds of $C_t(w)$.
For each literal in that conjunction, either its feature is basic, in which case
agreement supplies the required truth value, or it is a positive situation literal,
which the preceding paragraph supplies. Thus its step-clause antecedent is true,
and the clause forces $\sigma(X_{t+1,\phi})$ true.
The one-hot argument upgrades these positive facts to full agreement at time $t+1$.
This is
\source{Complexity/Computability/Complexity/CookLevin/Correctness.lean}{\lean{Agrees.step}};
induction gives
\source{Complexity/Computability/Complexity/CookLevin/Correctness.lean}{\lean{agrees_of_eval_tableau}}.

The direction of the auxiliary constraints matters. They force the read symbols
and situations that actually occur, but need not rule out every extra auxiliary
truth value. The proof works because nonbasic literals in step terms are positive
situation literals. Agreement itself concerns basic features.
At time $T$, agreement turns the two final clauses into the actual halted state
and output $[\mathrm{true}]$. This proves the forward direction of~\eqref{eq:tableau}.

\subsection{From an accepting run to a satisfying assignment}

Conversely, suppose $w$ has length $L$ and $C_T(w)$ is halted with output
$[\mathrm{true}]$. Set each certificate variable to its bit in $w$, and set
$X_{t,\phi}$ true exactly when $\phi$ holds of $C_t(w)$. This is
\source{Complexity/Computability/Complexity/CookLevin/Correctness.lean}{\lean{runAssignment}}.
Its reconstructed certificate is $w$, by equality of the individual bits.

We check each clause family. The initial configuration supplies the initial
clauses. Existence and uniqueness of the feature holding in each basic group give
the one-hot clauses. A head reads the symbol at its position, which supplies the
read clauses; the definition of a situation supplies the situation clauses.
For a step clause, if its antecedent conjunction is true, the reverse direction
of the local step-equation lemma says its consequent feature holds at the next
time. Finally, the assumed halted state and output supply the final clauses.
These checks are
\source{Complexity/Computability/Complexity/CookLevin/Correctness.lean}{\lean{eval_tableau_runAssignment}}.

The two directions first prove an equivalence with the certificate fixed
(\lean{exists_eval_tableau_iff}). Existentially quantifying that certificate gives
\source{Complexity/Computability/Complexity/CookLevin/Correctness.lean}{\lean{tableau_satisfiable_iff}},
which is exactly~\eqref{eq:tableau}.

\subsection{Instantiating the tableau with an NP verifier}

Let $A$ be any problem in NP under an encoding $\mathrm{enc}$.
Unpack its witness characterization: there are a polynomial $p$, a Boolean
verifier $V$, and a tidy polynomial-time machine computing $V$, such that
$a\in A$ exactly when some witness of length $p(|\mathrm{enc}(a)|)$ is accepted.
The completeness proof extracts fixed naturals $c,d$ that bound the verifier's
time by $c(m+1)^d$ on encoded pairs of length $m$.

Put $n=|\mathrm{enc}(a)|$ and $L=p(n)$.
The pair encoding frames the instance and appends the witness, so its length is
$m=2n+1+L$. Choose
\[
 u=\mathrm{frame}(\mathrm{enc}(a)),\qquad
 T=c(2n+1+p(n)+1)^d.
\]
Construct $F(M,u,L,T)$ and relabel its variables injectively by natural numbers.
Injectivity preserves satisfiability: an assignment pulls back along the numbering,
and any assignment to the original variables extends to their image and then to
all numbered variables. The library applies \lean{CNF.satisfiable_relabel} here.
By~\eqref{eq:tableau}, the resulting formula is satisfiable exactly when some
length-$L$ witness makes the machine halt with output true. The verifier certificate
identifies its output with $[V(a,w)]$, so the witness condition is exactly $a\in A$.
This proves
\source{Complexity/Computability/Complexity/CookLevin/Reduction.lean}{\lean{reduce_mem_SAT_iff_of_computesTidily}}.

\subsection{Certifying the reduction's resource bound}

The efficiency proof follows the same construction. First compute $(u,L,T)$ from
$a$, representing the two numeric bounds in unary. Framing, encoded length,
polynomial evaluation and the fixed power expression have polynomial-time
certificates, composed in
\source{Complexity/Computability/Complexity/CookLevin/Parameters.lean}{\lean{computableTidilyInPolyTime_parameters}}.

Next construct the tableau. Its clause constructors use bounded lists of times,
positions and features. Each clause-family constructor has a certificate in
\source{Complexity/Computability/Complexity/CookLevin/Clauses.lean}{\lean{CookLevin/Clauses.lean}}.
Unfolding \lean{tableau} expresses it as the concatenation of these families and
bounded list traversals. The \lean{polytime} tactic combines their registered
certificates with the list-operation rules; this proves
\source{Complexity/Computability/Complexity/CookLevin/PolyTime.lean}{\lean{computableTidilyInPolyTime_tableau}}.
The final certificate composes parameter computation, tableau construction and
binary variable relabelling, in that order. Thus it accounts for constructing and
encoding the reduction's output.

The machine, tape count, finite state space and polynomial bounds are fixed in
this argument. Its finite transition and enumeration choices may therefore be
stored in fixed tables. This explains why Lean definitions marked
\lean{noncomputable} still have certificates asserting the existence of machines
with polynomial bounds for all encoded instances. The proof does not extract a
particular runnable implementation of those choices.

\subsection{The final completeness step}

For arbitrary $A\in\cc{NP}$, the preceding correctness equivalence and efficiency
certificate supply a polynomial-time many-one reduction to SAT. This is
\source{Complexity/Computability/Complexity/CookLevin/Completeness.lean}{\lean{SAT_npHard}}.
The separately certified SAT verifier gives \lean{SAT.mem_NP}. Pairing these two
facts proves \lean{SAT_npComplete}. The generic complementation theorem then
turns SAT completeness into \lean{UNSAT_coNPComplete}.

